\subsection{An executed arithmetic probe certificate through height twenty}
\label{sec:executed-probe}

The local estimate can be executed with arithmetic separation data.
Take $N=4$, so $V_4$ has support $[-4,4]$, mesh $h=1/16$, and
dimension $126$.  The probe and its transform are
\[
 \phi(x)=\frac{e^{-|x|}-e^{|x|-2}}2\,\mathbf1_{|x|\le1},
 \qquad
 p(t)=\frac{1+e^{-2}-2e^{-1}\cos t}{1+t^2}.
\]
All the quantities in the following calculation are obtained from
finite arithmetic Weil pairings and completed moments at zero.

First, the completed moments themselves give an adequate local gap.
List the distinct positive eigenvalues of $C$ as
$\lambda_1>\lambda_2>\cdots>0$, and write
$S_k/2=\sum_jm_j\lambda_j^{2k}$.  For a positive Rayleigh value
$\mu$, one has $\lambda_1\ge\mu$.  The level-four calculation gives
\[
 b_7:=\bigl(S_7/2-\mu^{14}\bigr)^{1/14}
       <0.048230126,
 \qquad \mu-b_7>0.022515334.
\]
For every $j\ge2$,
\[
 \lambda_j^{14}\le S_7/2-\lambda_1^{14}
                        \le S_7/2-\mu^{14},
\]
so $\lambda_j\le b_7$.  The negative eigenvalues and zero have
distance at least $\lambda_1$ from $\lambda_1$.  Thus $d=9/400$
is a certified lower bound for the full spectral gap at $\lambda_1$.
Moreover, $\mu>1/20>b_7$ proves that the packet of positive ordinates
through height $20$ contains exactly one distinct point.
This moment inequality also implies $m_1=1$: the contrary inequality
$m_1\ge2$ would give $S_7/2\ge2\mu^{14}$, contradicting $b_7<\mu$.
The additional probe computation below verifies the projection-mass
rounding mechanism on these same arithmetic data.

\begin{proposition}[A completed local probe calculation]
\label{prop:executed-local-probe}
For a certified unit vector $v\in V_4$, the quantities in
Proposition~\ref{prop:local-probe-residual} and
\eqref{eq:residual-arithmetic} have the following enclosures.
\begin{center}
\begin{tabular}{ll}
\hline
Quantity & Certified enclosure or bound\\
\hline
$q(\phi)$ & $(9.98158838,9.98158839)\,10^{-5}$\\
$\mu$ & $(0.0707454607,0.0707454609)$\\
$\rho$ & $(0.0003983098,0.0003983100)$\\
$d$ & $9/400$ is a strict lower bound\\
$\delta_v$ & $<0.018022$\\
$\xi_v$ & $<0.001435$\\
$u/p(1/\mu)^2$ & $(0.9996487647,0.9996487649)$\\
Total multiplicity error & $<0.080978$\\
\hline
\end{tabular}
\end{center}
The resulting interval contains exactly the integer one.  This
executes the spectral-probe test for every positive ordinate through
height $20$.
\end{proposition}

\begin{proof}
We describe the arithmetic test functions, the certified linear
algebra, and the rounding inequality.  The primitive Gram matrix
$J_N$ is evaluated by the degree-five test preceding
\eqref{eq:two-tent-bound}.  For the additional probe pairings, put
$r=e^{-1}$, $g(x)=e^{-|x|}/2$, and $T_sg(x)=g(x+s)$.  Direct
substitution on the intervals separated by $-1,0,1$ gives
\[
 \phi=\bigl((1+r^2)I-r(T_1+T_{-1})\bigr)g.
\]
Splitting its convolution integral at $0$ and $x$ gives
$g*g=k$, where $k(x)=(1+|x|)e^{-|x|}/4$.  Consequently
\begin{align*}
 \phi*\phi={}&(1+4r^2+r^4)k
   -2r(1+r^2)(T_1+T_{-1})k\\
   &\quad+r^2(T_2+T_{-2})k.
\end{align*}
This is the compact test for $q(\phi)$.

For the central tent $t_h=(1-|x|/h)_+$, take the causal primitive
$A_\phi(x)=\int_{-\infty}^x(x-u)\phi(u)\,du$.  The distributional
identity $t_h''=(\delta_{-h}-2\delta_0+\delta_h)/h$ proves
\[
 K(x):=(\phi*t_h)(x)
   =\frac{A_\phi(x+h)-2A_\phi(x)+A_\phi(x-h)}h.
\]
The evenized test for $q(\phi,t_h(\cdot-jh)-t_h)$ is therefore
$K(x+jh)+K(x-jh)-2K(x)$.  Every piece of these tests is a sum
$F(x)=\sum_{a=-1}^1(c_a+d_ax)e^{ax}$, with exact dyadic
breakpoints.  Their supports have radius at most $5$; the test for
$q(\phi)$ has radius $2$, and the primitive Gram tests have radius
at most $8$.  Thus every prime sum is finite.

For completeness, the removable singularity in the first
archimedean integration interval is cancelled by the exact identity
\[
 \frac{F(x)-F(0)}x
   =\sum_{a=-1}^1
        \bigl(c_a a\operatorname{exprel}(ax)+d_a e^{ax}\bigr),
 \qquad
 \operatorname{exprel}(z)={}_1F_1(1;2;z).
\]
The constant multiplying the original test value $F(0)/2$ is
\[
 C_h=-\log\pi-\gamma_E
   +\int_0^h\frac{2(e^{-2x}-e^{-x/2})}{1-e^{-2x}}\,dx
   -\log(1-e^{-2h}).
\]
After this constant is extracted, the first-cell remainder is
\[
 -\int_0^h
    \frac{e^{-x/2}}{2\operatorname{exprel}(-2x)}
                      \frac{F(x)-F(0)}x\,dx.
\]
The other cells have no singular endpoint.  Polynomial-exponential
antiderivatives evaluate the pole terms.  Certified complex ball
integration evaluates the remaining meromorphic callbacks, and
each prime logarithm is separated from the piece boundaries before
its corresponding expression is evaluated.

A midpoint congruence preconditions the positive Gram matrix; the
full outward congruence residual certifies positivity.  An approximate
eigenvector then proposes coefficients, which are frozen as exact
dyadic numbers before certification.  Writing these coefficients
$x=a+ib$ and $H_N=-iB$ with real antisymmetric $B$, the exact formulas
used are
\[
 \begin{split}
 n&=a^TGa+b^TGb,\qquad
 \mu=2a^TBb/n,\\
 \rho^2&=(a^TJ_Na+b^TJ_Nb)/n-\mu^2,\qquad
 u=((a^Tc)^2+(b^Tc)^2)/n.
 \end{split}
\]
All these expressions are evaluated with outward balls at $120$
decimal digits, and repeated at $160$ digits.  The finite input
Gram and primitive pairings, candidate coefficients, pairing
enclosures, and strict tests are recorded with source hashes in the
external certificates.

The residual interval lies entirely above $b_7$, so its spectral
point is $\lambda_1$.  The bounds above give $\rho<d/2$.
Proposition~\ref{prop:local-probe-residual} therefore applies.
The exact evaluated denominator $p(1/\mu)^2$ can be used in its
division step.  Since
\[
 m_1\le\frac{q(\phi)e^{2C_0\xi_v}}{p(1/\mu)^2}<4,
\]
integrality permits the a priori upper bound $m_1\le3$ even without
using the direct moment implication $m_1=1$.  The final error bound
computed in balls is
\[
 \frac{\delta_vq(\phi)}{p(1/\mu)^2}
             +3\bigl(e^{2C_0\xi_v}-1\bigr)<0.080978.
\]
Together with the enclosure for the divided mass, this lies
strictly inside $(1/2,3/2)$ and proves the claimed rounding.
The arithmetic packet count already proved that it covers all
positive ordinates through height $20$.
\end{proof}

At level three the same bound for the logarithmic error exceeds
$0.755$, so the rounding interval does not isolate an integer.  No probe certificate for the full packet through height
$45$ is asserted by the present computation.
